\begin{abstract}
최근의 시각 기반 웹 에이전트는 화면 이해, 요소 그라운딩, 행동
궤적을 서로 다른 데이터 혼합과 학습 단계로 습득한다. 기존
레시피를 검토하는 과정에서 우리는 행동 학습 이전의 웹페이지
설명이 모델마다 서로 다른 상세도로 구성되지만, 이 선택의 효과는
다른 학습 요인과 분리되어 평가되지 않았다는 점에 주목하였다. 본
연구는 같은 웹 화면을 핵심 정보만 남긴 Short description으로
학습하는 것과 화면의 텍스트ㆍ상태ㆍ행동 요소를 포괄한 Long
description으로 학습하는 것을 비교한다. 150K개 화면에서 공유된
Stage 1 사실을 추출한 뒤, Stage 2에서 평균 110단어의 Short 설명과
평균 3,656자의 Long 설명을 생성한다. 각 조건에는 동일한 1.2M
그라운딩 예제를 혼합한다. 후속 단계에서는 동일한 100K 명령어
데이터, 2K action adaptation 데이터, 100K 행동 궤적을 순차적으로
학습한다.
PT 직후, 명령어 튜닝 이후, action adaptation 이후, 행동 궤적 학습
이후의 체크포인트를 함께 평가함으로써 서술 상세도의 효과가
어디에서 나타나고 얼마나 유지되는지를 측정한다. 데이터 검증
과정에서는 프롬프트의 금지 규칙에도 teacher가 약 3.2\%의 예제에서
존재하지 않는 좌표를 생성함을 관찰했으며, 좌표 집합 대조를 통해
이를 제거하였다. 이 통제 설계는 웹 GUI 사전학습에서 더 많은
내용을 서술하는 것이 실제로 더 유용한 supervision인지 검증한다.
\end{abstract}
